\section{总结与延伸}

\subsection{核心结论}

\begin{enumerate}
\item 谢赛宁的个人路线不是成功学，而是研究问题如何通过选择、组织和偶然逐渐成形。
\item Vision 不应被理解成一组旧任务，而应被理解成处理连续、高维、有噪声真实世界信号的 perspective。
\item 世界模型的最小定义是 $s_{t+1}=f(s_t,a_t)$，但真正困难的是 state 表征、物理接地、planning 和反馈闭环。
\item LLM 是智能系统的重要组成部分，但它主要吃下人类已经 tokenized 的知识，不能自动替代物理世界建模。
\item “反向 OpenAI”意味着世界模型不能只下载互联网，而要和真实世界伙伴共建数据、评测和应用闭环。
\item AMI Labs 的关键不是“又一家大模型公司”，而是尝试在研究自由、资源规模、商业闭环和学术连接之间找新组织形态。
\item LLM-pilled 的问题是单一叙事会定义榜单和资源配置，从而压缩世界模型、视频理解、物理智能等方向的探索空间。
\item 智能不是一条语言分数线；真实世界智能要回到身体、环境、行动、反馈和目标。
\end{enumerate}

\subsection{开放问题}

总结之后还需要保留问题意识，因为本期最诚实的地方之一，就是承认世界模型的许多关键答案还没有收敛。下面这些问题，也是后续阅读和继续生成笔记时应该持续追踪的线索。

\begin{itemize}
\item 世界模型预训练是否会形成类似 next-token prediction 的稳定范式？
\item 世界模型的 state 应该是显式 3D、latent representation、混合结构，还是按任务动态形成？
\item VLA 与 world model 的边界会逐渐融合，还是形成上游基础层与下游行动层的分工？
\item AMI 式的全球伙伴网络能否真正带来不可替代的数据闭环？
\item 当 LLM 榜单逐渐饱和后，产业会不会把资源重新投向 physical intelligence？
\end{itemize}

\subsection{拓展阅读}

\begin{itemize}
\item Yann LeCun, \textit{A Path Towards Autonomous Machine Intelligence}：理解 JEPA、世界模型和 autonomous intelligence 的长期路线。
\item Richard Sutton, Dyna / model-based reinforcement learning 相关论文：理解世界模型、planning 和 RL 的历史关系。
\item Model Predictive Control 入门材料：理解“用模型 roll out 未来并选择动作”的控制直觉。
\item 视觉表征学习、视频生成、3D representation、VLA 和机器人预训练相关论文：理解世界模型可能的多条入口。
\item Frans de Waal, \textit{Are We Smart Enough to Know How Smart Animals Are?}：理解反人类中心的智能观。
\end{itemize}

\begin{importantbox}{最后的判断}
本期最值得带走的不是“LLM 错了、世界模型对了”这种二分，而是一个更精确的判断：LLM 解决了人类已写下世界的一大部分；世界模型要解决的是人类没有写下、也很难写下的世界。下一阶段智能的关键，可能就在这两者之间的桥上。
\end{importantbox}

\end{document}
